
\section{Bounding sums}\doHeader{Bounding sums}

Recall that, for example, $\sum_{i=1}^n i^2$ denotes the sum $1^2 + 2^2 + \cdots + n^2$.
You may have previously learned \emph{closed forms} for some sums,
such as $\sum_{i=1}^n i = n(n+1)/2$.
Unfortunately, for most sums of interest, closed forms
are difficult to find or simply don't exist.
Fortunately, when we are analyzing asymptotic running times,
we don't care about lower-order terms or constant factors,
so we don't usually need exact closed forms for the sums that arise.
Instead, for example, we only need to know that $\sum_{i=1}^n i = \Theta(n^2)$.
Here are three rules for quickly estimating the big-$\Theta$ values of sums.

\subsection{Naive upper and lower bounds}
Here are two rules that work well for sums whose terms grow polynomially (or less) with the index.

\begin{observation}[Naive upper bound]
  Any sum is at most the number of terms times the max.~term.
\end{observation}
\paragraph{Three examples:}
\begin{align*}
  \textstyle \sum_{i=1}^n i ~\le~ {}
  & n \cdot \max_{i=1}^n i ~=~ n\cdot n ~=~ O(n^2)
  \\
  \textstyle \sum_{i=1}^n i^2 ~\le~ {}
  & n \cdot \max_{i=1}^n i^2 ~=~ n\cdot n^2 ~=~ O(n^3)
  \\
  \textstyle \sum_{i=1}^n \log i ~\le~ {}
  & n \cdot \max_{i=1}^n \log i ~=~ n \log n ~=~ O(n \log n)
\end{align*}

\begin{observation}[Naive lower bound]
  Any sum with non-negative terms
  is at least half the number of terms times the median term.
\end{observation}
\paragraph{Three examples:}
\begin{align*}
  \textstyle \sum_{i=1}^n i ~\ge~ {}
  & (n/2) \cdot (n/2) ~=~ \Omega(n^2)
  \\
  \textstyle \sum_{i=1}^n i^2 ~\ge~ {}
  & (n/2) \cdot (n/2)^2 ~=~ \Omega(n^3)
  \\
  \textstyle \sum_{i=1}^n \log i ~\ge~ {}
  & (n/2) \cdot \log (n/2) ~=~ \Omega(n \log n)
\end{align*}
 
For the three examples above, the naive upper and lower bounds match (up to constant factors).
So together they are enough to determine the value of the sum, up to constant factors (the big-$\Theta$ value).
For example, $\sum_{i=1}^n \log i = \Theta(n\log n)$.

The bounds will match for any sum whose terms don't increase (or decrease) too fast---specifically, for any sum where the median term is within a constant factor
of the maximum term.
For example, it will work for sums whose terms grow only polynomially with the index.


\subsection{Geometric sums}
However, it won't work for all sums.
For example, try $\sum_{i=1}^n 2^i$.
The naive upper bound gives $\sum_{i=1}^n 2^i = O(n 2^n)$.
The naive lower bound gives $\sum_{i=1}^n 2^i = \Omega(n 2^{n/2})$.
But these upper and lower bounds are far apart!
Indeed their ratio $2^n / 2^{n/2} = 2^{n/2}$ is quite large.
The answer is actually $\sum_{i=1}^n 2^i = \Theta(2^n)$.
Neither the naive upper or lower bound is tight in this case.

We have another rule for such sums.  First we 
formally define the type of sums it applies to:
\begin{definition}
  A sum $\sum_{i=1}^n \alpha_i$ is \emph{geometric} if,
  for some constants $c>1$ and $i_0$ (indep.~of $n$)
  either\\ (i) for each $i\ge i_0$ the $i$th term
  is at least $c$ times the previous term
  (that is, $\alpha_{i} \ge c \cdot \alpha_{i-1}$), or \\
  (ii) for each $i\ge i_0$ the $i$th term
  is at most the previous term divided by $c$
  (that is, $\alpha_{i} \le \alpha_{i-1}/c$).
\end{definition}

\emph{Rule of thumb:} if the sum has an exponential term in the summand, it is probably geometric.

\paragraph{Three examples}
\begin{itemize}
\item $\sum_{i=1}^n 2^i$ is geometric:
  the ratio of successive terms $2^{i}/ 2^{i-1}$ is at least two.
\item $\sum_{i=1}^n 2^i/i$ is also geometric:
  for $i\ge 3$ the ratio of successive terms is $(2^i/i)/(2^{i-1}/(i-1))$ $=2(i-1)/i\ge 4/3$.
\item $\sum_{i=1}^n i^2+1$ is not geometric:
  the ratio of successive terms $((i+1)^2+1) / (i^2 +1)$
  gets arbitrarily close to 1 for large $i$.
\end{itemize}


Here's our third rule.  It handles any geometric sum:

\begin{observation}[geometric sums]
  Any geometric sum is proportional to its largest term.
\end{observation}
\paragraph{Three examples}~{}

\[
\begin{array}{rr@{~}l}
  & \sum_{i=0}^n 2^i &= \Theta(\max_{i=0}^n 2^i) = \Theta(2^n). \\[2ex]
  & \sum_{i=1}^n 2^i/i &= \Theta(\max_{i=1}^n 2^i/i) = \Theta(2^n/n).   \\[2ex]
  & \sum_{i=0}^\infty i^2/2^i&  = \Theta(\max_{i=0}^\infty i^2/2^i) = \Theta(1).
\end{array}
\]

\subsection{Exercises}

\exercise \emph{(Bounding sums)}
Use the so-called naive upper and lower bounds, and the rule for geometric sums,
to determine the big-$\Theta$ value of each of the following sums.
% You don't have to show your work,
% but you can if you want feedback on it
% (and it may help you get partial credit if your answer is wrong). 


\begin{enumerate}[label=(\alph*)]
\item $\displaystyle\sum_{i=1}^n 1.1^i/i^2$

\item $\displaystyle\sum_{i=1}^{\log_2 n} i^2\log_3 i$

\item $\displaystyle\sum_{i=1}^{2^n} (\log_3 i)^2$
  
\item $\displaystyle\sum_{i=0}^{n} i^3/2^i$
\end{enumerate}

\subsection{External sources on bounding sums}
\begin{mylist}
\item \CLRS Appendices A-1 and A-2
\end{mylist}
  
\newpage

\section{Recursion trees for recursive algorithms}\doHeader{Recursion trees}\label{dc: sec: recursion trees}
Next we discuss how to analyze the run times of recursive algorithms,
using \textsf{mergesort} (shown in Figure~\ref{dc: fig: mergesort}) as an example.
Let $T(n)$ denote the worst-case running time of \textsf{mergesort} on any array of size $n$.
Then $T(1) = \Theta(1)$.
% [review 2026-09-27] fix: the recurrence holds for n >= 2 (n = 1 is the base case T(1) = Theta(1) stated just above); was n >= 1
For $n\ge 2$, by inspection of \textsf{mergesort} and
using that \textsf{merge} takes linear time, we have the recurrence relation
\[
  T(n) = 2\, T(n/2) + n.
\]
(For now, we assume $n$ is a power of 2, so $n/2$ is an integer.)
\begin{figure}
\begin{myframed}
\begin{steps}
  \item[] {\sf mergesort}$(A[1..n])$
  \step if $n = 1$, return
  \step set $m= \lfloor n/2\rfloor$
  \step \textsf{mergesort}$(A[1..m])$ \comment{(sort the first $n/2$ elements)}
  \step \textsf{mergesort}$(A[m+1..n])$ \comment{(sort the last $n/2$ elements)}
  \step \textsf{merge}$(A[1..n], m)$ \comment{(merge the two sorted halves, in time $\Theta(n)$)}
\end{steps}
\end{myframed}
\caption{Mergesort. The \textsf{merge} procedure runs in linear time.  We don't give its code.}\label{dc: fig: mergesort}
\end{figure}

\paragraph{Using a \emph{recursion tree} to count the work.}
A \emph{recursion tree} represents the recursive structure of an execution of a recursive algorithm.
The root node represents the top-level call.
Each non-root node represents one recursive call.
Each node is labeled by the size of the subarray that the corresponding call sorts.
The children of each node for a given call
represent the two recursive calls that that call makes.
Figure~\ref{dc: fig: recursion tree} shows the recursion tree for \textsf{mergesort} on an array of size $n=16$.
\begin{figure}
\begin{myframed}
\tikzset{
    vertex/.style={inner sep=0pt, minimum size=17pt},
    edge/.style={> = latex'}
}
\noindent\adjustbox{valign=c, raise=-5pt}{
\begin{tikzpicture}[scale=0.5]
  \node[vertex] (1) at (7.5,4) {16} ;

  \node[vertex] (2) at (3.5,3) {8} ;
  \node[vertex] (3) at (11.5,3) {8} ;
  \draw[edge] (2) -- (1); 
  \draw[edge] (3) -- (1); 

  \node[vertex] (4) at (1.5,2) {4} ;
  \node[vertex] (5) at (5.5,2) {4} ;
  \node[vertex] (6) at (9.5,2) {4} ;
  \node[vertex] (7) at (13.5,2) {4} ;
  \draw[edge] (4) -- (2); 
  \draw[edge] (5) -- (2); 
  \draw[edge] (6) -- (3); 
  \draw[edge] (7) -- (3); 

  \node[vertex] (8) at (0.5,1) {2} ;
  \node[vertex] (9) at (2.5,1) {2} ;
  \node[vertex] (10) at (4.5,1) {2} ;
  \node[vertex] (11) at (6.5,1) {2} ;
  \node[vertex] (12) at (8.5,1) {2} ; 
  \node[vertex] (13) at (10.5,1) {2} ; 
  \node[vertex] (14) at (12.5,1) {2} ; 
  \node[vertex] (15) at (14.5,1) {2} ;
  \draw[edge] (8) -- (4);
  \draw[edge] (9) -- (4);
  \draw[edge] (10) -- (5);
  \draw[edge] (11) -- (5);
  \draw[edge] (12) -- (6);
  \draw[edge] (13) -- (6);
  \draw[edge] (14) -- (7);
  \draw[edge] (15) -- (7);
  
  \node[vertex] (16) at (0,0) {1} ;
  \node[vertex] (17) at (1,0) {1} ;
  \node[vertex] (18) at (2,0) {1} ;
  \node[vertex] (19) at (3,0) {1} ;
  \node[vertex] (20) at (4,0) {1} ;
  \node[vertex] (21) at (5,0) {1} ;
  \node[vertex] (22) at (6,0) {1} ;
  \node[vertex] (23) at (7,0) {1} ;
  \node[vertex] (24) at (8,0) {1} ;
  \node[vertex] (25) at (9,0) {1} ;
  \node[vertex] (26) at (10,0) {1} ;
  \node[vertex] (27) at (11,0) {1} ; 
  \node[vertex] (28) at (12,0) {1} ; 
  \node[vertex] (29) at (13,0) {1} ; 
  \node[vertex] (30) at (14,0) {1} ; 
  \node[vertex] (31) at (15,0) {1} ;
  \draw[edge] (16) -- (8);
  \draw[edge] (17) -- (8);
  \draw[edge] (18) -- (9);
  \draw[edge] (19) -- (9);
  \draw[edge] (20) -- (10);
  \draw[edge] (21) -- (10);
  \draw[edge] (22) -- (11);
  \draw[edge] (23) -- (11);
  \draw[edge] (24) -- (12);
  \draw[edge] (25) -- (12);
  \draw[edge] (26) -- (13);
  \draw[edge] (27) -- (13);
  \draw[edge] (28) -- (14);
  \draw[edge] (29) -- (14);
  \draw[edge] (30) -- (15);
  \draw[edge] (31) -- (15);
  % \draw[edge] (a) -- (b); 
\end{tikzpicture}
}
\adjustbox{valign=c}{\small
  \begin{tabular}{@{}c|ccc|c@{}}
    & \# of & subprob. & work per & work in \\
    level & subprobs & size  & subprob & level \\ \hline 
    0 & 1 & 16 & 16 & 16 \\[2pt]
    1 & 2 & 8 & 8 & 16 \\[2pt]
    2 & 4 & 4 & 4 & 16 \\[2pt]
    3 & 8 & 2 & 2 & 16 \\[2pt]
    4 & 16 & 1 & 1 & 16 \\ \hline
    \multicolumn{4}{r}{total:} & $5\times 16$
  \end{tabular}
}
\end{myframed}
\caption{On the left, the recursion tree for \textsf{mergesort} on an array of size 16.
  On the right,  a table of the work done in each level of the recursion tree.}\label{dc: fig: recursion tree}
\end{figure}

As illustrated in the table to the right in the figure, if the input array size $n$ is a power of 2,
there are $1+\log_2 n$ levels.
For each $i\in\{0,\ldots, \log_2 n\}$,
level $i$ has $2^i$ subproblems, each of size $n/2^i$.
The work associated with a subproblem of size $n/2^i$ is $\Theta(n/2^i)$,
so the total work at level $i$ is $\Theta(2^i \times n/2^i) = \Theta(n)$.
Summing over the $\log n$ levels, the total work is $\Theta(n\log n)$.
Hence, $T(n) = \Theta(n\log n)$.
Note that the work is balanced evenly across the levels.

If $n$ is not a power of 2, then there are $1+\lceil \log_2 n \rceil$ levels.
In each level $i$ except the last, there are $2^i$ subproblems,
% [review 2026-09-27] fix: subproblem sizes can exceed n/2^i because of ceilings; replaced "between n/2^{i+1} and n/2^i" by the correct bound "at most ceil(n/2^i) <= 2n/2^i"
each of size at most $\lceil n/2^i\rceil \le 2n/2^i$.
So, by a similar calculation, the total work $T(n)$ is also $\Theta(n\log n)$ in this case.

\begin{theorem}
  \textsf{Mergesort} runs in $O(n \log n)$ time.
\end{theorem}

\paragraph{Variation with heavy leaves.}

Consider the recurrence relation $T(n) = \mathbf{3}\cdot T(n/2) + n$.
The same approach applies, but now each non-leaf node has \textbf{three} children.
As before we have
\begin{align*}
  T(n) & {}= \sum_{i=0}^{\text{max level}} (\text{\# subprobs on level } i) \times (\text{work per subprob.~on level } i)
\end{align*}
but now the number of subproblems on level $i$ is $3^i$, so the calculation proceeds as follows:
\begin{align*}
  T(n) & {}= \sum_{i=0}^{\log_2 n} 3^i (n/2^i) ~=~ n \sum_{i=0}^{\log_2 n}\, (3/2)^i
    = \Theta(n\, (3/2)^{\log_2 n})  \text{~~~~(because the sum is geometric)} \\
  & {}= \Theta(n\cdot n^{\log_2 3/2}) = \Theta(n^{\log_2 3}) = \Theta(n^{1.58\ldots})
    \text{~~~~(using $b^{\log n} = n^{\log b}$ and $\log_2 3/2 = \log_2(3) -1$).}
\end{align*}

For \textsf{mergesort}, the work was balanced across the levels.
But here we get a qualitatively different behavior:
\emph{the total work is dominated by the work done at the leaves.}

\paragraph{Variation with heavy root.}
Consider the recurrence relation $T(n) = 3\, T(n/2) + \mathbf{n^2}$.
Then
\begin{align*}
  T(n) & {} = \sum_{i=0}^{\log_2 n} 3^i (n/2^i)^{\mathbf 2} = n^2 \sum_{i=0}^{\log_2 n}\, (3/4)^i
           = \Theta(n^2).
\end{align*}
In this case \emph{the total work is dominated by the work at the root.}

\medskip

These are the three common cases:
either (a) the work is balanced across the levels,
or the sum is geometric, so either (b) the work at the leaves dominates,
or (c) the work at the root dominates.
A particularly important special case of the latter case
(when the work at the root dominates)
is when \emph{the algorithm does linear work,
then recurses on subproblems whose combined size
is a constant factor less than the size of the given problem}.

For example,
consider the recurrence relation $T(n) = a T(n/b) + n$,
for some $a < b$.  That is, the algorithm does linear work,
then recurses on $a$ problems, each of size $n/b$,
so having total size $an/b = (a/b) n$ where $a/b < 1$.
Then we get
\begin{align*}
  T(n) & {} = \sum_{i=0}^{\log_b n} a^i (n/b^i) = n \sum_{i=0}^{\log_b n}\, (a/b)^i
           = \Theta(n).
\end{align*}
The last step follows because $a < b$, so the sum is geometric
and dominated by the first term.
See also the Selection algorithm in Section~\ref{dc: sec: selection}.

% Here's a more general case, which arises in the algorithm for \prob{Selection}.
% \begin{lemma} Suppose $T(1) = 1$ and for $n\ge 2$
%   \[\displaystyle T(n) \le n + \sum_{i=1}^k a_i\, T(n/b_i) \]
%   where $a_i > 0$ and $b_i > 1$ for each $i$, and $\sum_{i=1}^k a_i/b_i < 1$.
%   Then $T(n) = O(n)$.
% \end{lemma}
% \begin{proof}[Proof sketch.]
%   We sketch a proof, by induction on $n$, that $T(n) \le c\, n$ for some constant $c$.
%   We choose (to make the proof below work) $c = 1+1/(1- \sum_{i=1}^k a_i/b_i )$.
%   The base case $n=1$ holds because $c\ge 1$.
%   For the inductive step, consider $n\ge 2$, then
%   \begin{align*}
%     T(n)
%     & \le  n +\textstyle\sum_{i=1}^k a_i\, T(n/b_i)   && \comment{(by assumption)} \\
%     & \le n + \textstyle\sum_{i=1}^k a_i\, c\, n/b_i   && \comment{(by induction, using $b_i>1$)} \\
%     & = c\, n (1/c +\textstyle\sum_{i=1}^k a_i/b_i) \\
%     & \le c\, n && \comment{(by the choice of $c$)}.
%   \end{align*}
%   The last step because by the choice of $c$,
%   $c \ge 1/(1- \sum_{i=1}^k a_i/b_i )$,
%   so $1/c + \sum_{i=1}^k a_i/b_i \le 1$.

%   This is only a proof sketch because it ignores the details
%   re WORKING
% \end{proof}

\pythonCode{divide_and_conquer/mergesort.py}

\subsection{Exercises}

\paragraph{With solution}

\exercise\label{dc: exercise: recursion tree with solution}\emph{(Recursion-tree method)}
For the algorithm below,
use the recursion-tree method to find the asymptotic value (using the $\Theta$-notation)
for the number of letters printed by the algorithm as a function of $n$.

\newcommand{\ALG}[4]{
  \vbox{                        % keep on one page
    \begin{steps}
    \item[] def \textsf{Print{#1}s}$(n)$:  \comment{Assume $n$ is a power of $#4$.}
      \step for $i = 1$ to $#2$: print(``{#1}'') \comment{\ldots print $#2$ #1's}
      \step if $n > 1$:
      \step~~~for $j=1$ to $#3$: \textsf{Print{#1}s}$(\floor{n/#4})$
      \comment{\ldots make $#3$ recursive calls, each with input $\floor{n/#4}$}
    \end{steps}
  }
}

\begin{myframed}
\ALG{Y}{n\log_2 n}{3}{3}
\end{myframed}

\newpage

The solution should give
\begin{enumerate}[label=(\roman*)]
\item an appropriate recurrence relation,  
\item as a function of the level $i$ in the recursion tree:
the number of nodes (that is, calls) in level $i$,
the problem size (the argument to each call) at each node in level $i$,
and the number of letters printed by each call in level $i$,
\item the number of levels in the recursion tree,
\item the sum, over all levels, of the number of letters printed by calls in that level,
  and finally,
\item a big-$\Theta$ estimate of the value of the sum
(as a function of $n$), obtained by using either the rule for geometric
sums or the naive upper and lower bounds
(show your work for (v) if you want feedback).
\end{enumerate}

\solution

\begin{enumerate}[label=(\roman*)]
\item Define $Y(n)$ to be the number of letters printed by {\sc PrintYs} given input $n\ge 1$.

  Then $Y(1) = 0$ and, for $n>1$, $Y(n) = n \log_2 n + 3 Y(n/3)$.
\item In level $i$, there are $3^i$ nodes.
  Each call has input $n/3^i$ and prints $(n/3^i)\log_2(n/3^i)$ letters.
% [review 2026-09-27] fix: the levels are 0, ..., log_3 n (as in the sum in (iv) and the bound in (v)), so there are 1 + log_3 n levels (was log_3 n)
\item There are $1 + \log_3 n$ levels.
\item The total number of letters printed is 
\( Y(n)
~=~ 
\sum_{i=0}^{\log_3 n} 3^i (n/3^i)\log_2 (n/3^i)
~=~
\sum_{i=0}^{\log_3 n} n\log_2(n/3^i)\).
\item By the naive upper bound, the sum is at most $(1+\log_3 n) n\log_2 n = O(n\log^2 n)$.

By the naive lower bound, the sum is at least $\frac{1}{2}(1+\log_3 n)$ times the median term,
which is $n\log_2(n/3^i)$ where $i = (1+\log_3 n)/2$.  For this $i$, $3^i$ simplifies to $\Theta(\sqrt n)$,
so the median term simplifies to $n\log_2 (n/\Theta(\sqrt n)) = n\log_2 \Theta(\sqrt n) = \Theta(n\log n)$.
So the sum is $\Omega(n\log^2 n)$.

Thus, the sum, which equals $Y(n)$, is $\Theta(n\log^2 n)$.
\end{enumerate}

\paragraph{Exercise without solutions}

\exercise\label{dc: exercise: mergesort}\emph{(Mergesort variation)} % CS 5800 HWK 7
Consider a variant of Mergesort that, given an array of size $N$,
recurses as usual just until the subproblems are of size $\sqrt N$,
then solves each subproblem of size $\sqrt N$ directly using insertion sort
(a quadratic-time algorithm).
Use a recursion-tree analysis to determine
the asymptotic worst-case running time of this algorithm on arrays of size $n$.
% [review 2026-09-27] fix: cross-reference pointed to Exercise 4.4 (dc: exercise: recursion), which only refers back; the work to show, (i)-(v), is described in Exercise 4.2 (dc: exercise: recursion tree with solution)
Show your work as described for Exercise~\ref{dc: exercise: recursion tree with solution}.

\exercise\label{dc: exercise: recursion}\emph{(Recursion-tree method)}
Repeat Exercise~\ref{dc: exercise: recursion tree with solution} for each of the algorithms below.
Give your answers in the same format as the solution given for that exercise.
To make the calculations easier, ignore floors in your calculations.
(That is, for (a), (c), and (e) assume $n$ is a power of 3.
For (b) assume $n$ is a power of 4. 
For (d) assume $n$ is a power of 2.)

\begin{enumerate}[label=(\alph*)]
\item \ALG{A}{3n}{4}{3}
  
\item \ALG{B}{n^2}{16}{4}

\item \ALG{C}{5n}{2}{3}

\item \ALG{D}{11}{2}{2}

\item \ALG{E}{2n^3}{9}{3}

\end{enumerate}

%%%%%%%%%%%%%%%% PROBLEM %%%%%%%%%%%%%%%%%

\subsection{External sources on recursion trees and recursive algorithms}\label{dc: sec: refs}

\begin{itemize}

\item \DPV Chapter 2.2; \KT Chapter 5; \CLRS Chapter 4 (especially 4.4); \LLM Chapter 22;

\item \JE Chapter 1, and Appendix II Section 3 (\& exercises starting on Page 20), at
  
  {\small \url{https://jeffe.cs.illinois.edu/teaching/algorithms/notes/99-recurrences.pdf}}

\item An MIT lecture video
  

\URL{https://ocw.mit.edu/courses/electrical-engineering-and-computer-science/6-006-introduction-to-algorithms-fall-2011/lecture-videos/lecture-3-insertion-sort-merge-sort/}

\end{itemize}

%%% Local Variables:
%%% mode: latex
%%% TeX-master: "divide_and_conquer"
%%% End:
